\begin{figure}[ht]
\begin{AcademicBox}[案例 1：政策决策（长上下文）]
\small
\setlength{\parindent}{0pt}

% --- Context ---
\textbf{\textit{上下文：}} \\
“地图阴影区域表示 ESCAP 成员与准成员。*\
亚洲及太平洋经济社会委员会（ESCAP）是亚太地区最具包容性的政府间平台。
该委员会推动 53 个成员国和 9 个准成员合作，共同应对可持续发展挑战。
ESCAP 是联合国五个区域委员会之一。
ESCAP 秘书处通过产出行动导向知识、提供技术援助与能力建设服务，支持各国发展目标、区域协议以及《2030 年可持续发展议程》的落实。
*地图上的名称与材料展示方式，不代表联合国秘书处对任何国家或地区法律地位的立场…… ”
\vspace{4pt} \hrule \vspace{4pt}

% --- Question ---
\textbf{\textit{问题：}} \\
在保持政治与社会稳定的前提下，政府应采取何种政策组合，才能更好平衡财政可持续性、私营部门参与以及能源转型？ \\
\vspace{4pt} \hrule \vspace{4pt}

% --- Correct Result (Full Width) ---
\textbf{\textcolor{teal}{正确预测（我们的方法 / 标准答案）：}} \hfill \textcolor{teal}{\checkmark} \\
\textbf{“渐进式能源转型 + 国家绿色投资银行”方案：}
建立国家绿色投资银行……\textbf{在未来五年维持现有化石燃料补贴}，以保障能源价格稳定并逐步扩大可再生能源……\textbf{同时推迟征收碳税}。
\vspace{4pt} \hrule \vspace{4pt}

% --- Incorrect Results (Grouped) ---
\textbf{\textcolor{red}{错误基线：}} \\
\footnotesize
\begin{itemize}[leftmargin=*, label={}, itemsep=4pt, topsep=2pt]
    \item \textbf{Qwen3-8B、PruLong、DuoAttn、InfLLM-V2：} \\
    \textbf{“激进财政再平衡”方案：} 提高石油生产碳税……\textbf{在五年内逐步取消化石燃料补贴}……
    \textit{（错误：过早征收碳税和取消补贴，与稳定目标冲突）}

    \item \textbf{MoBA：} \\
    \textbf{“立即取消化石燃料补贴”方案：} \textbf{立刻取消所有化石燃料补贴}……并在两年内引入碳定价机制……
    \textit{（错误：“立即取消”忽视社会稳定约束）}

    \item \textbf{NSA：} \\
    \textbf{“混合融资 + 出口导向”方案：} 建立公私混合融资基金……\textbf{实施适度碳税}……
    \textit{（错误：错误强调出口市场，且碳税时机不当）}

\end{itemize}

\end{AcademicBox}
\caption{复杂政策推理任务上的定性对比。我们的模型正确识别出满足稳定性约束所需的“渐进式”路径，而基线往往幻觉出“激进”或“立即”措施，违背题设稳定要求。}
\label{fig:case_policy}
\end{figure}




\begin{figure}[ht]
\begin{AcademicBox}[案例 2：法律推理与文档理解]
\small
\setlength{\parindent}{0pt}

% --- Context Section ---
\textbf{\textit{上下文：}} \\
截至 2024 年 6 月 20 日的加拿大法律文本（含英法双语条文）片段：
包括《财政管理法》（Financial Administration Act）及其法条编纂信息、修订状态、官方法律效力说明等内容；文本中穿插英法双语段落与法律术语，结构冗长且噪声较高。
\vspace{4pt} \hrule \vspace{4pt}

% --- Question Section ---
\textbf{\textit{问题：}} \\
“特别经济措施法（SEMA）”与“财政管理法（FAA）”在职责上不同……在一个加拿大公民企业资产被扣押的假设场景中，它们的适用差异应如何理解？哪一种解释最准确地体现两者之间的细微关系？
\vspace{4pt} \hrule \vspace{4pt}

% --- Correct Result ---
\textbf{\textcolor{teal}{正确预测（我们的方法 / 标准答案）：}} \hfill \textcolor{teal}{\checkmark} \\
SEMA 允许政府在更广泛的制裁框架下先行冻结资产，而不立即涉及没收归属；
相对地，FAA 需要法律程序裁定资产是否应被\textbf{永久没收并重新分配用于公共用途}。
\vspace{4pt} \hrule \vspace{4pt}

% --- Incorrect Results ---
\textbf{\textcolor{red}{错误基线：}} \\
\footnotesize
\begin{itemize}[leftmargin=*, label={}, itemsep=4pt, topsep=2pt]

    \item \textbf{Qwen3-8B、PruLong、DuoAttn、MoBA、NSA：} \\
    回答仅指出 SEMA 可冻结资产、FAA 需\textbf{评估法律依据}并确保\textbf{符合加拿大法律}……
    \textit{（错误：回答过于泛化，未识别“公共用途再分配”这一关键结论）}

    \item \textbf{InfLLM-V2：} \\
    强调 FAA 需法律审查并可能用于政府用途，SEMA 可先行冻结并\textbf{延后即时决定}……
    \textit{（错误：过度聚焦流程时延，而非资产去向这一实体差异）}

\end{itemize}

\end{AcademicBox}
\caption{双语法律文档上的对比。我们的模型准确抽取了 FAA 关于资产\textbf{再分配用于公共用途}的关键法律含义，而基线多停留在“法律评估/合规”层面的泛化表述。}
\label{fig:case_legal}
\end{figure}
